\section{Global level sets and the control partition}\label{sec:global}

Let $J=K+1-k_0$ be the number of blocks and
\[
 \mathcal L(R)=\{a:\Qctrl(a)\le R\}.
\]
We encode each assignment in $\mathcal L(R)$ by four finite layers:
\begin{enumerate}[label=(\arabic*)]
\item the hot-block set $H$ and its local energy-fibre code; the hot charge
satisfies
\[
 \sum_{k\in H}R_w(c_2,k)\le R;
\]
\item the cold boundaries $B$ with unequal adjacent labels, for which
$\sum_{k\in B}\Pi_k\le R$;
\item the integer shells $v_k=\lfloor Q_k(a)\rfloor$;
\item the label and exception data on maximal cold segments.
\end{enumerate}
The first three layers identify the cold segments.  On each segment the
exception-aware boundary theorem propagates one dominant label.  The first
label lies in an interval of length
$O(1+\sqrt R/\sctrl)$; later segment starts are charged to a preceding hot
block or mismatch boundary.  Exceptional primes and their residues are paid
for by \eqref{eq:exception-count}.

The decoder must be uniform in every integer label recorded by the code.  A
non-wrapped label outside the quantitative cold-label range has empty fibre,
since any member would contradict \eqref{eq:label-bound}.  A wrapped large
label represents the same residues as a small dominant label $M$.  Two primes
from the agreement class rigidify the lift: if two admissible small integers
agree modulo both primes, their difference is divisible by their product but
smaller in absolute value, hence zero.  Retaining the assignment and replacing
the recorded label by $M$ injects the wrapped fibre into one fixed small-label
fibre.  Thus no uncounted large-label family remains.

